\documentclass[11pt,letterpaper,sans]{moderncv}

% moderncv styles: casual, classic, banking, oldstyle, fancy
\moderncvstyle{classic}
\moderncvcolor{blue}

% Page geometry
\usepackage[scale=0.9]{geometry}

% Encoding / typography
\usepackage[T1]{fontenc}
\usepackage{ragged2e}
\usepackage{etoolbox}

% Reduce blocky justification in bullet lists
\AtBeginEnvironment{itemize}{\RaggedRight}
\hyphenpenalty=7000
\exhyphenpenalty=7000
\righthyphenmin=4
\emergencystretch=2em

% Personal data
\name{Peter}{Selby}
\address{8329 129th Pl SE}{Newcastle, WA}{}
\phone[mobile]{(805)~604~5298}
\email{peter@selby.ca}

\begin{document}
\makecvtitle
\pagestyle{empty}

\section{Summary}
\cvitem{}{Software engineer with 10+ years of experience building and scaling distributed systems, developer platforms, and internal tools at Google, Apple, Amazon, and ByteDance. Skilled in backend architecture, data processing, and full-stack development. Passionate about building reliable, high-performance systems and tools.}

\section{Technical Skills}
\cvitem{Languages}{Java, C++, Go, Scala, Python, JavaScript, Clojure, Lisp, Ruby, SQL, Erlang}
\cvitem{Databases}{PostgreSQL, MySQL, Oracle SQL, Amazon Redshift}
\cvitem{Architecture}{Distributed systems, Kubernetes, REST APIs, data processing with Hadoop \& Spark, microservices, gRPC, Protocol Buffers}
\\cvitem{Infrastructure \\& Tooling}{Kubernetes, Docker, NixOS, AWS, Flux, ArgoCD, Prometheus, Grafana, CI/CD pipelines, Linux \\& Unix, encryption, profiling, database design \\& optimization, benchmarking}

\section{Experience}

\cventry
  {2023--2025}
  {Site Reliability Engineer}
  {ByteDance Inc}
  {Bellevue, WA}
  {}
  {
    \begin{itemize}
      \item Worked on real-time communication (RTC) infrastructure for
        large-scale live streaming.
      \item Designed and implemented a distributed rolling upgrade system for
        worldwide RTC media ingress nodes, enabling safe and automated fleet
        upgrades with minimal redundancy or capacity loss. Reduced deployment
        cycle from multiple weeks to a 24-hour release, increasing development
        velocity and service reliability.
      \item Created, modified and/or maintained various cloud services in Python
        and Go. Used ArgoCD for continuous deployment and GitOps-based infrastructure management.
      \item Managed oncall responsibilities for large-scale production services. Used Prometheus and Grafana extensively for problem diagnosis, incident response, and service troubleshooting to minimize downtime and improve reliability.
    \end{itemize}
  }

\cventry
  {2016--2021}
  {Site Reliability Engineer}
  {Google Inc}
  {Kirkland, WA}
  {}
  {
    \begin{itemize}
      \item Integrated complex service configurations with Google's internal
        configuration management platform, rewriting multiple RTC services
        (including telephony, chat, and video) to safely consume, validate,
        apply, and monitor live configuration updates. Supported canaries, A/B
        testing, and automated rollbacks. Partnered with 5+ development teams to
        standardize configuration logic and ensure consistent behavior across
        heterogeneous systems serving 100M+ users.
      \item Redesigned and replaced the cloud rollout system for 20+ RTC
        services, implementing multi-stage canarying, key-metric monitoring, and
        automated regression analysis, enabling data-driven rollout and rollback
        decisions. Reduced rollout risk and improved release velocity across a
        multi-team production environment.
      \item Collaborated with cross-functional teams to refactor core RTC
        services to improve performance, scalability, and observability. Worked extensively with gRPC for service communication.
    \end{itemize}
  }

\cventry
  {2015--2016}
  {Full Stack Engineer}
  {Apple Inc}
  {Cupertino, CA (Remote)}
  {}
  {
    \begin{itemize}
    \item Built internal web applications to support Apple Maps Transit, including an infinite-scrolling POI data editor, enabling faster data-entry workflows to accelerate the launch of Apple Transit.
    \item Designed features from the backend to the web UI, using a combination of PostgreSQL, Scala and JavaScript.
    \item Supported new features, batch changes, and optimized cronjobs using PostgreSQL. Designed data models and validation workflows enabling rapid editing and publishing of curated transit mapping data provided by on-site contractors. Optimized database queries and schemas for complex transit data relationships.
    \end{itemize}
  }

\cventry
  {2013--2015}
  {SDE II / Technical Program Manager}
  {Amazon.com}
  {Seattle, WA}
  {}
  {
    \begin{itemize}
      \item Led the design and development of a large-scale financial reporting platform, generating daily company-wide fine-grained P\&L reports.
      \item Worked with users and stakeholders to gather requirements, pitch alternatives, and determine direction.
      \item Designed and built ETL data ingestion and aggregation services in Scala and SQL, among the first at Amazon to leverage AWS Redshift. Worked with Ruby and Clojure for various internal tools and infrastructure.
      \item Bootstrapped a new engineering team, assisted with roadmap and goals, and took a leading role in architecture and technology-stack design.
      \item Developed an interactive query service in Scala to enable richer interaction with financial data, allowing users without a SQL background to perform deeper analysis of daily reporting data.
      \item Launched AWS cloud-native platform service for company-wide use within Amazon.
    \end{itemize}
  }

\cventry
  {2011--2013}
  {Software Development Engineer I \& II}
  {Amazon.com}
  {Seattle, WA}
  {}
  {
    \begin{itemize}
    \item Engineered distributed ETL pipelines (Hadoop and SQL) producing daily financial and operational metrics at a per-shipment level across Amazon's global business.
    \item Designed and implemented algorithms to allocate and reconcile diverse costs and revenue components, including liquidations, returns, shipping, and promotions, to individual shipments. Used Oracle SQL and later Amazon Redshift for large-scale data processing.
    \item Partnered with finance and operations stakeholders to define heuristics, validate financial correctness, and integrate results into a company-wide reporting system used throughout Amazon to drive day-to-day business decisions.
    \item Extended data coverage to new geographies and business domains, including FBA and 3P. Launched the data-processing pipeline for Amazon China.
    \end{itemize}
  }

\section{Education}

\cventry
  {2008--2011}
  {Master of Science, Computer Science}
  {University of British Columbia}
  {Vancouver, BC}
  {}
  {
    Focus areas included Distributed Systems, Operating Systems, Cloud Computing, Multi-Agent Systems, and Game Theory. Selected projects included building two variants of a peer-to-peer distributed filesystem (directory-server and DHT-based) exploring metadata scalability, and implementing a proof-of-concept alternative to object-oriented programming in Lisp where objects emerged dynamically via pattern matching.
  }

\cventry
  {2003--2008}
  {Bachelor of Science, Honours, Computer Science}
  {University of Winnipeg}
  {Winnipeg, MB}
  {}
  {}

\end{document}